\chapter*{Аннотация}\label{abstract}

\begin{center}
	\textbf{Проходимые мосты Эйнштейна-Розена внутри черных дыр}
	
	\textit{Прокопьев Константин Эдуардович}
\end{center}

В рамках настоящего исследования проведено расширение концепции моста Эйнштейна-Розена для более общих пространственно-временных метрик, включая решение Райсснера-Нордстрёма (описывающее заряженные невращающиеся чёрные дыры) и метрику Керра-Ньюмена (характеризующую заряженные вращающиеся чёрные дыры). Особое внимание уделено анализу геодезических траекторий в указанных пространствах-временах, на основании которого сделано предположение о потенциальной проходимости обобщённого моста Эйнштейна-Розена. Отдельно исследован феномен инфляции массы — эффекта неограниченного роста массы чёрной дыры при малых возмущениях вблизи горизонта Коши, характерного для внутренней структуры заряженных и вращающихся решений. В работе предложены возможные пути устранения данной проблемы в рамках классической ОТО без привлечения квантовых эффектов.

\textbf{Ключевые слова:} ОТО, черные дыры, мост Эйнштейна-Розена, инфляция массы.

\newpage